\documentclass[10pt,a4paper]{amsart}
\usepackage[margin=19mm]{geometry}
\usepackage{amsmath,amssymb,mathtools}
\usepackage[scr=rsfs,cal=euler]{mathalfa}
\usepackage{xfrac}
\usepackage[unicode,hidelinks,bookmarks=false]{hyperref}
\newcommand{\ie}{i.e.\ }
\newcommand{\cf}{cf.\ }
\newcommand{\resp}{respectively\ }
\newcommand{\Subsection}{\subsection}
\providecommand{\nobreakdash}{\nobreak\hbox{-}}
\setlength{\emergencystretch}{2em}
\multlinegap0pt
\usepackage{bm}
\usepackage{bbm}
\usepackage{dsfont}
\usepackage{xspace}
\usepackage{cancel}
\usepackage{mathtools}
\usepackage{amsthm}
\makeatletter
\@ifundefined{theorem}{\newtheorem{theorem}{Theorem}}{}
\makeatother
\DeclareMathOperator{\cmobmv}{Mob^*_{\mu}}
\DeclareMathOperator{\cp}{\,\square\,}
\DeclareMathOperator{\mobmv}{Mob_{\mu}}
\newcommand{\move}{\rightsquigarrow}
\title[Solution to some conjectures on mobile position problems]{Solution to some conjectures on mobile position problems}
\author{Ethan Shallcross, James Tuite, Aoise Evans, Aditi Krishnakumar, Sumaiyah Boshar}
\date{Source: arXiv:2507.16622v1 (CC BY 4.0)}
\begin{document}
\maketitle

\noindent\emph{Source and licence.} Solution to some conjectures on mobile position problems. Version arXiv:2507.16622v1 (CC BY 4.0), distributed under the Creative Commons Attribution 4.0 International licence, as recorded in the arXiv licence metadata for this exact version and shown on the abstract page https://arxiv.org/abs/2507.16622v1. Full rights notice: licenses/arxiv-2507.16622-CC-BY-4.0.txt.\par\smallskip

\noindent\textit{Editorial scope.} Counted unit: the Theorem whose source label is \texttt{smallcartgrids} in the article (source0.tex lines 758--773, source label \texttt{smallcartgrids}), delivered together with the article's own complete original proof of it (source0.tex lines 774--796, one \texttt{proof} environment of 23 lines). No mathematics is written, repaired or re-derived: statement and proof are the article's own text line for line. Required context retained uncounted: largecartgrids (lines 497--505). Every definition, display and label of those lines is retained unchanged. Omitted from this package and not redistributed: 1--496, 506--757, 797--1123. Nothing inside a retained excerpt is omitted or altered: every retained body is delivered exactly as the article's lines are written, so no omission receipt of any kind accompanies this package. Citations: the retained excerpts carry no citation command at all, so no bibliography entry of the article is transcribed and no reference is redistributed. Reference adaptation: none was needed and none was made; every reference inside the delivered unit resolves inside it, so no unresolved reference or citation is produced. Printed slips kept literal: no printed word, symbol, hypothesis or number of the counted statement or of its proof was altered. Layout and class: the article's own class is \texttt{article} with packages including amssymb, amsthm, tikz, url, pgfplots, amsmath, todonotes, subcaption, color, graphicx, placeins, caption, geometry; the wrapper is the lane's pinned \texttt{amsart} editorial wrapper at 10pt on A4, which restates the theorem-like environments the retained text needs and adds the article's own macro definitions that the retained text needs (\texttt{\textbackslash cmobmv}, \texttt{\textbackslash cp}, \texttt{\textbackslash mobmv}, \texttt{\textbackslash move}), with extra wrapper packages (bm, bbm, dsfont, xspace, cancel, mathtools). The delivered numbering is the unit's own. Assets: no retained span contains a figure environment, a graphics-inclusion command, a PDF-page inclusion, a table or a TikZ picture; no image file is copied into this package, the package declares no asset dependency and no image file is redistributed with it. Licence: version arXiv:2507.16622v1 of this article is distributed under the Creative Commons Attribution 4.0 International licence, as recorded in the arXiv licence metadata for this exact version and shown on the abstract page https://arxiv.org/abs/2507.16622v1. A full rights notice is delivered at licenses/arxiv-2507.16622-CC-BY-4.0.txt. Characters: every retained excerpt is pure ASCII, so the delivered engine, XeLaTeX with its default Latin Modern text font, reports no missing character.
	\begin{theorem}
		For $n \geq m \geq 5$, \[\mobmv(P_n \cp P_m) =
		\begin{cases}
			2m, & \text{if}\ n > m, \\
			2m-1, & \text{if}\ n = m.
		\end{cases}
		\]
		\label{largecartgrids}
	\end{theorem}

	\begin{theorem}
		For $a \geq 3$,		 
		\[\mobmv(P_n \cp P_3) =
		\begin{cases}
			4, & \text{if}\ n = 3, \\
			5, & \text{if}\ n = 4, \\
			6, & \text{if}\ n \geq 6.
		\end{cases} \] and \[ \mobmv(P_n \cp P_4) =
		\begin{cases}
			5, & \text{if}\ n = 3, \\
			7, & \text{if}\ n = 4 \,\text{or}\ 5 \\
			8, & \text{if}\ n \geq 6.
		\end{cases}
		\]
		\label{smallcartgrids}
	\end{theorem}

	\begin{proof}
		To prove the upper bounds, we find the cardinalities of largest mutual visibility sets containing a central vertex of the grid.
		
		Firstly, suppose we had a mutual visibility set $S$ of $P_3 \cp P_3$ with $|S| = 5$, $(2,2) \in S$. We can suppose that $2$ vertices in $S$ have second coordinate $1$. In order to maintain mutual visibility, $(2,2)$ can be the only vertex in $S$ with second coordinate $2$. So there are $2$ vertices in $S$ with second coordinate $3$. It is easily seen that these vertices cannot be mutually visible and so $\mobmv(P_3 \cp P_3) \leq 4$.
		
		To see that $\mobmv(P_3 \cp P_3) = 4$, station robots at $(1,2), (2,1), (2,3), (3,2)$ and move $(2,1) \move (3,1)$, $(2,3) \move (1,1)$, and then $(1,2) \move (2,2)$. Since the positions of the robots may be permuted by moving them around the edge of the grid, in fact $\cmobmv(P_3 \cp P_3) = 4$.
		
		Next, suppose we had a mutual visibility set $S$ of $P_4 \cp P_3$ with $|S| = 6$, $(2,2) \in S$. Exactly two vertices in $S$ have each possible second coordinate. By the above reasoning, there can be at most four vertices in $S$ with first coordinate $<4$, so exactly two vertices in $S$ have first coordinate $4$. Similarly if $(3,2) \in S$, there are exactly two vertices in $S$ with first coordinate $1$ and the vertices cannot be mutually visible. If $(1,2) \in S$, there can only be at most one in $S$ with second coordinate $1$, a contradiction. Finally, if $(4,2) \in S$ we can assume without loss of generality that $(4,1) \in S$ and $(4,3) \notin S$. But then $(4,1)$ and the vertex in $S$ with second coordinate $3$ and smallest first coordinate will not be mutually visible. Thus, there can be no such $S$ and $\mobmv(P_4 \cp P_3) \leq 5$.
		
		To establish equality, station robots at $(1,2), (1,3), (2,1), (4,2), (4,3)$ and move $(1,2) \move (2,2)$, $(2,1) \move (1,1)$, and then $(1,3) \move (2,3)$. Since the positions of the robots may be permuted by moving them around the edge of the grid, in fact $\cmobmv(P_4 \cp P_3) = 4$.
		
		It also follows that any mutual visibility set of $P_4 \cp P_4$ containing $(1,1)$ can have at most five vertices with first coordinate $<4$ and so $\mobmv(P_4 \cp P_4) \leq 7$.
		
		To see that $\mobmv(P_4 \cp P_4) = 7$, station robots at the vertices \[ (1,2), (1,3), (2,4), (3,1), (3,4), (4,2), (4,3)\] and move $(1,2) \move (2,2)$, $(2,4) \move (1,4)$. From the initial configuration we may move $(3,1) \move (2,1)$ or $(1,2), (1,1), (2,1)$ and deduce the result by symmetry. We can use moves similar to those in Theorem~\ref{largecartgrids} to permute the robots, so that $\cmobmv(P_4 \cp P_4) = 7$.
		
		Now suppose we had a mutual visibility set $S$ of $P_5 \cp P_4$ with $|S| = 8$, $(3,2) \in S$. Exactly two vertices in $S$ have each possible second coordinate. If $(2,2) \in S$, we also have $(1,1), (1,3), (5,1), (5,3) \in S$. But then $S$ cannot contain two vertices with second coordinate $4$, a contradiction. If $(1,2) \in S$ then as two vertices in $S$ have second coordinate $1$, $(1,1) \in S$. But similar reasoning gives $(1,3) \in S$, a contradiction. Hence, $\mobmv(P_4 \cp P_5) \leq 7$. Equality follows from the result for $P_4 \cp P_4$ by applying the `shifting' argument used in the proof of Theorem~\ref{largecartgrids}. We also see that $\cmobmv(P_5 \cp P_4) = 7$.
		
		For $n \geq 2$, since at most two vertices in any mutual visibility set can have the same second coordinate, it is clear that $\mobmv(P_n \cp P_3) \leq 6$ and $\mobmv(P_n \cp P_4) \leq 8$.
		
		To see that $\mobmv(P_5 \cp P_3) = 6$, position six robots as in the proof of Theorem~\ref{largecartgrids}. Then perform the moves $(3,1) \move (4,1) \move (5,1)$, $(3,3) \move (4,3) \move (5,3)$, $(2,1) \move (1,1)$, $(1,2) \move (2,2)$, $(2,3) \move (1,3)$ and then $(2,2) \move (3,2)$. Using similar moves similar to those in Theorem~\ref{largecartgrids} to permute the robots, we see that $\cmobmv(P_5 \cp P_3) = 5$.
		
		Finally, we show that $\mobmv(P_6 \cp P_4) = 8$. Position eight robots as in the proof of Theorem~\ref{largecartgrids}. We may follow the sequences in that proof to traverse all vertices except $(3,2), (3,3)$, and $[1] \times [4]$. However, by moving the robots to the right in decreasing order of the first coordinates, we see that the remaining vertices can be traversed. Moves similar to those in Theorem~\ref{largecartgrids} show that $\mobmv(P_6 \cp P_4) = 7$ and $\mobmv(P_5 \cp P_3) = 5$.
	\end{proof}

\end{document}
